% ---
% titre: 9p - Nombres naturels et décimaux - Corrigés (réponses seules)
% theme: NO
% chapitre: Nombres naturels et décimaux
% sous_chapitre: 
% niveau: [9e]
% type: RES
% tags: c-criteres-divisibilite, c-decomposition-facteurs-premiers, c-nombres-premiers, c-ppmc, c-puissances, c-priorite-operations, f-corrige, p-auto-correction
% difficulte: 1
% corrige: false
% source: latex
% ---

% ---- Macros locales à ce fichier (pas dans template_site.tex) ----
\colorlet{pathGreen}{green!15}

\newcommand{\etape}[1]{%
  \vspace{10pt}
  {\Large\bfseries\textcolor{themeColor}{#1}}
  \vspace{4pt}
  \par
}

\newcommand{\pc}[2]{\cellcolor{themeColor!15}\shortstack{\textbf{#1}\\\textbf{#2}}}
\newcommand{\pcg}[2]{\cellcolor{pathGreen}\shortstack{\textbf{#1}\\\textbf{#2}}}
\newcommand{\ncell}[2]{\shortstack{\textcolor{gray}{#1}\\\textcolor{gray}{#2}}}
\newcommand{\entree}{\cellcolor{themeColor!50}\textbf{4\cdot5}}
\newcommand{\sortie}{\cellcolor{themeColor!50}\textbf{27}}

% =========================================================
\etape{1 \quad Critères de divisibilité}
% =========================================================

\header{NO3 --- Comment s'en sortir?}

{\footnotesize
\renewcommand{\arraystretch}{2.5}
\noindent
\begin{tabularx}{\textwidth}{|*{14}{C|}}
\hline
\ncell{54}{9\cdot7} & \ncell{36}{12\cdot3} & \ncell{22}{5\cdot8} & \entree & \ncell{18}{6\cdot7} & \ncell{14}{2\cdot8} & \ncell{88}{4\cdot9} & \ncell{70}{3\cdot7} & \ncell{48}{5\cdot6} & \ncell{96}{4\cdot9} & \ncell{27}{6\cdot2} & \ncell{54}{4\cdot12} & \ncell{36}{10\cdot11} & \ncell{108}{7\cdot11} \\
\hline
\ncell{72}{2\cdot8} & \ncell{63}{8\cdot9} & \pc{20}{7\cdot8} & \ncell{30}{3\cdot9} & \ncell{24}{3\cdot8} & \ncell{12}{6\cdot4} & \ncell{110}{9\cdot8} & \pcg{72}{8\cdot8} & \pcg{64}{6\cdot4} & \pcg{24}{3\cdot5} & \ncell{18}{8\cdot9} & \ncell{50}{8\cdot7} & \ncell{77}{10\cdot5} & \ncell{15}{3\cdot8} \\
\hline
\ncell{16}{7\cdot4} & \ncell{28}{8\cdot7} & \pc{56}{8\cdot4} & \ncell{24}{3\cdot9} & \pc{27}{8\cdot6} & \pc{48}{5\cdot7} & \ncell{44}{4\cdot10} & \pcg{40}{6\cdot12} & \ncell{28}{5\cdot3} & \pcg{15}{12\cdot3} & \ncell{72}{8\cdot7} & \ncell{56}{8\cdot9} & \ncell{24}{3\cdot11} & \ncell{33}{8\cdot3} \\
\hline
\ncell{21}{11\cdot6} & \ncell{66}{7\cdot3} & \ncell{36}{4\cdot8} & \pc{32}{3\cdot9} & \pcg{48}{3\cdot10} & \pcg{30}{7\cdot5} & \pc{35}{9\cdot6} & \ncell{144}{3\cdot8} & \pcg{24}{5\cdot8} & \pcg{36}{9\cdot7} & \ncell{56}{9\cdot11} & \ncell{99}{7\cdot8} & \ncell{132}{2\cdot11} & \ncell{22}{11\cdot12} \\
\hline
\ncell{49}{8\cdot7} & \ncell{42}{7\cdot9} & \ncell{28}{4\cdot5} & \ncell{25}{3\cdot5} & \ncell{14}{8\cdot6} & \ncell{66}{9\cdot6} & \pc{54}{7\cdot7} & \pc{49}{2\cdot12} & \ncell{121}{7\cdot7} & \pcg{63}{8\cdot11} & \ncell{77}{9\cdot6} & \ncell{54}{11\cdot0} & \ncell{42}{2\cdot8} & \ncell{16}{7\cdot6} \\
\hline
\ncell{32}{8\cdot8} & \ncell{64}{9\cdot9} & \ncell{81}{2\cdot10} & \ncell{20}{4\cdot9} & \pc{36}{6\cdot5} & \pc{84}{9\cdot4} & \ncell{60}{9\cdot5} & \pc{24}{11\cdot8} & \pc{88}{7\cdot3} & \ncell{56}{7\cdot4} & \ncell{84}{9\cdot10} & \ncell{90}{8\cdot12} & \pc{81}{12\cdot8} & \pc{96}{6\cdot6} \\
\hline
\ncell{45}{5\cdot5} & \ncell{25}{9\cdot5} & \pc{18}{12\cdot4} & \pc{30}{6\cdot3} & \ncell{32}{7\cdot10} & \ncell{70}{8\cdot4} & \pc{66}{12\cdot7} & \pc{45}{6\cdot11} & \pc{21}{9\cdot5} & \pc{108}{4\cdot7} & \pc{28}{11\cdot10} & \pc{110}{9\cdot9} & \ncell{72}{8\cdot8} & \pc{36}{8\cdot10} \\
\hline
\ncell{16}{7\cdot7} & \ncell{49}{8\cdot2} & \pc{48}{9\cdot2} & \ncell{12}{7\cdot6} & \pc{42}{8\cdot9} & \pc{72}{5\cdot5} & \ncell{24}{11\cdot3} & \ncell{33}{5\cdot12} & \pc{14}{12\cdot9} & \ncell{55}{8\cdot4} & \ncell{32}{10\cdot6} & \pc{60}{11\cdot12} & \pc{80}{5\cdot12} & \ncell{64}{7\cdot8} \\
\hline
\ncell{24}{2\cdot11} & \ncell{22}{4\cdot6} & \ncell{36}{2\cdot9} & \pc{18}{12\cdot2} & \pc{24}{7\cdot6} & \ncell{60}{5\cdot5} & \pc{25}{4\cdot11} & \pc{44}{10\cdot6} & \pc{60}{2\cdot7} & \ncell{12}{6\cdot10} & \pc{132}{6\cdot11} & \ncell{48}{9\cdot4} & \ncell{32}{6\cdot8} & \ncell{56}{4\cdot8} \\
\hline
\ncell{56}{3\cdot10} & \ncell{30}{9\cdot3} & \ncell{27}{10\cdot2} & \ncell{20}{7\cdot8} & \ncell{56}{8\cdot9} & \ncell{48}{8\cdot6} & \pc{120}{8\cdot7} & \ncell{110}{10\cdot2} & \ncell{15}{9\cdot4} & \ncell{36}{3\cdot5} & \pc{66}{5\cdot6} & \ncell{45}{6\cdot11} & \ncell{8}{9\cdot5} & \ncell{63}{2\cdot4} \\
\hline
\ncell{96}{4\cdot8} & \ncell{72}{9\cdot12} & \ncell{21}{7\cdot8} & \ncell{16}{9\cdot9} & \ncell{96}{7\cdot10} & \pc{56}{6\cdot12} & \ncell{44}{5\cdot11} & \pc{16}{10\cdot12} & \pc{24}{4\cdot4} & \pc{30}{8\cdot3} & \pc{54}{3\cdot3} & \pc{9}{2\cdot8} & \pc{16}{7\cdot11} & \ncell{18}{4\cdot5} \\
\hline
\ncell{32}{8\cdot12} & \ncell{84}{7\cdot8} & \ncell{56}{9\cdot9} & \ncell{81}{3\cdot5} & \pc{72}{3\cdot4} & \ncell{70}{4\cdot6} & \ncell{55}{7\cdot10} & \pc{32}{8\cdot9} & \pc{72}{12\cdot12} & \pc{144}{6\cdot9} & \ncell{36}{4\cdot6} & \ncell{18}{3\cdot9} & \pc{77}{9\cdot12} & \pc{108}{4\cdot7} \\
\hline
\ncell{27}{11\cdot10} & \ncell{110}{12\cdot11} & \ncell{121}{2\cdot5} & \ncell{10}{8\cdot10} & \pc{12}{8\cdot11} & \pc{88}{7\cdot3} & \pc{21}{4\cdot8} & \ncell{96}{11\cdot10} & \ncell{110}{8\cdot12} & \ncell{132}{11\cdot4} & \ncell{44}{9\cdot10} & \ncell{96}{8\cdot9} & \ncell{72}{2\cdot9} & \pc{28}{10\cdot7} \\
\hline
\ncell{36}{7\cdot6} & \ncell{42}{3\cdot9} & \ncell{50}{6\cdot4} & \ncell{24}{10\cdot8} & \ncell{80}{11\cdot8} & \ncell{28}{4\cdot4} & \ncell{48}{9\cdot4} & \ncell{36}{6\cdot7} & \ncell{42}{12\cdot8} & \ncell{56}{10\cdot9} & \ncell{90}{8\cdot7} & \sortie & \pc{70}{3\cdot9} & \ncell{35}{4\cdot8} \\
\hline
\end{tabularx}}

\vspace{20pt}
\header{Problème 3 --- Divisibilité}
{\small
\textbf{A.} \hfill \textbf{B.}\hfill\,

\begin{tabularx}{0.45\textwidth}{l|*{8}{C}}
 & 2 & 3 & 4 & 5 & 6 & 9 & 10 & 25 \\
\hline
23\,462  & \texttimes &   &   &   &   &   &   &   \\
21\,300  & \texttimes & \texttimes & \texttimes & \texttimes & \texttimes &   & \texttimes & \texttimes \\
1\,881   &   & \texttimes &   &   &   & \texttimes &   &   \\
35\,325  &   & \texttimes &   & \texttimes &   & \texttimes &   & \texttimes \\
21\,180  & \texttimes & \texttimes & \texttimes & \texttimes & \texttimes &   & \texttimes &   \\
72\,342  & \texttimes & \texttimes &   &   & \texttimes & \texttimes &   &   \\
320\,105 &   &   &   & \texttimes &   &   &   &   \\
\end{tabularx}\hfill
\begin{tabularx}{0.45\textwidth}{l|*{8}{C}}
 & 2 & 3 & 4 & 5 & 6 & 9 & 10 & 25 \\
\hline
37\,452  & \texttimes & \texttimes & \texttimes &   & \texttimes &   &   &   \\
61\,587  &   & \texttimes &   &   &   & \texttimes &   &   \\
3\,075   &   & \texttimes &   & \texttimes &   &   &   & \texttimes \\
108\,826 & \texttimes &   &   &   &   &   &   &   \\
31\,150  & \texttimes &   &   & \texttimes &   &   & \texttimes & \texttimes \\
1\,100   & \texttimes &   & \texttimes & \texttimes &   &   & \texttimes & \texttimes \\
729\,936 & \texttimes & \texttimes & \texttimes &   & \texttimes & \texttimes &   &   \\
\end{tabularx}
\hfill\,

\vspace{6pt}
\textbf{C.} \emph{(Trouve toutes les possibilités)}\\[2pt]
\begin{tabularx}{\textwidth}{l|*{8}{C}}
 & 3 & 4 & 5 & 6 & 9 & 25 \\
\hline
37\,452 & \texttimes & \texttimes &  & \texttimes &  &  \\
615\,807 -- 615\,843 -- 615\,861 -- 615\,879 & \texttimes &  &  &  & \texttimes &  \\
3075 -- 3225 -- 3525 -- 3675 -- 3975    & \texttimes &  & \texttimes &  &  & \texttimes  \\
18\,121 -- 18\,122 -- 18\,127           &  &  &  &  &  &  \\
11\,150 -- 31\,150 -- 41\,150 -- 61\,150 -- 71\,150 -- 91\,150 &  & & \texttimes  &  &  & \texttimes  \\
1770 & \texttimes &  & \texttimes & \texttimes &  &  \\
729\,936 & \texttimes & \texttimes &  & \texttimes & \texttimes &  \\
\end{tabularx}
}

\vspace{20pt}
\header{NO28 --- Quels chiffres\,?}
\vspace{-15pt}\begin{multicols}{4}
\begin{enumerate}[leftmargin=*,align=parleft, label=\alph*)]
\item 0\,;\,2\,;\,4\,;\,6 ou 8
\item 0\,;\,3\,;\,6 ou 9
\item 0 ou 5
\item 2\,;\,4\,;\,6 ou 8 
\item 3\,;\,6 ou 9
\item 0 
\item 2\,;\,4 ou 8 
\item 0 ou 6
\end{enumerate}
\end{multicols}

\vspace{30pt}
% =========================================================
\etape{2 \quad Nombres premiers et décomposition}
% =========================================================

\header{NO50 --- Décompositions}
\vspace{-10pt}\begin{tabularx}{\textwidth}{XXXX}
$12 = 2^2 \cdot 3$ &$84 = 2^2 \cdot 3 \cdot 7$&$600 = 2^3 \cdot 3 \cdot 5^2$ & $150\,000 = 2^4 \cdot 3 \cdot 5^5$\\
$49 = 7^2$&$180 = 2^2 \cdot 3^2 \cdot 5$ & $4200 = 2^3 \cdot 3 \cdot 5^2 \cdot 7$&\\
$54 = 2 \cdot 3^3$&$525 = 3 \cdot 5^2 \cdot 7$&$4700 = 2^2 \cdot 5^2 \cdot 47$ &\\
\end{tabularx}
 
\vspace{20pt}
\header{Problème 5 --- Factorisation}
\vspace{-10pt}\begin{minipage}[t]{0.32\textwidth}
\textbf{A.}\\[2pt]
$128 = 2^7$\\
$726 = 2 \cdot 3 \cdot 11^2$\\
$450 = 2 \cdot 3^2 \cdot 5^2$\\
$680 = 2^3 \cdot 5 \cdot 17$\\
$98 = 2 \cdot 7^2$\\
$195 = 3 \cdot 5 \cdot 13$
\end{minipage}%
\begin{minipage}[t]{0.34\textwidth}
\textbf{B.}\\[2pt]
$1125 = 3^2 \cdot 5^3$\\
$152 = 2^3 \cdot 19$\\
$2028 = 2^2 \cdot 3 \cdot 13^2$\\
$540 = 2^2 \cdot 3^3 \cdot 5$\\
$2310 = 2 \cdot 3 \cdot 5 \cdot 7 \cdot 11$\\
$81 = 3^4$
\end{minipage}%
\begin{minipage}[t]{0.32\textwidth}
\textbf{C.}\\[2pt]
$765 = 3^2 \cdot 5 \cdot 17$\\
$500 = 2^2 \cdot 5^3$\\
$2079 = 3^3 \cdot 7 \cdot 11$\\
$625 = 5^4$\\
$570 = 2 \cdot 3 \cdot 5 \cdot 19$\\
$260 = 2^2 \cdot 5 \cdot 13$
\end{minipage}

\vspace{30pt}
% =========================================================
\etape{3 \quad PPMC et PGDC}
% =========================================================

\header{NO52 --- ppmc}
\vspace{-15pt}\begin{multicols}{5}
\begin{enumerate}[leftmargin=*,align=parleft, label=\alph*)]
\item 6
\item 8
\item 18  
\item 12  
\item 72  
\item 100  
\item 70
\item 2\,500
\item 100
\item 60
\end{enumerate}
\end{multicols}

\newpage
\header{NO53 --- pgdc}
\vspace{-16pt}\begin{multicols}{5}
\begin{enumerate}[leftmargin=*,align=parleft, label=\alph*)]
\item 4
\item 3
\item 2 
\item 4
\item 1
\item 5 
\item 2
\item 2
\item 25
\item 120
\end{enumerate}
\end{multicols}

\vspace{-5pt}
\header{NO54 --- ppmc et pgdc}
\vspace{-16pt}\begin{multicols}{5}
\begin{enumerate}[leftmargin=*,align=parleft, label=\alph*)]
\item 1
\item 16
\item 17
\item 600 
\item 2
\item 23\,112
\item 18  
\item 1\,024  
\item 7 
\item 42
\end{enumerate}
\end{multicols}

\vspace{-5pt}
\header{Problème 10 --- C., H., J.}
\vspace{-20pt}\begin{itemize}[leftmargin=*,align=parleft]
\item[C.] Océane dispose de CHF 32,35 .
\item[H.] Une marche mesure 16 cm de haut 
\item[J.] Le premier bateau aura effectué 5 voyages, le deuxième 4 voyages et le troisième 3 voyages.
\end{itemize}

\header{Problème 10 --- Autres}
\vspace{-20pt}\begin{itemize}[leftmargin=*,align=parleft]
\item[A.] Loïse dipose de CHF 63.--.
\item[B.] Diane dispose de CHF 36.--.
\item[D.] Les carrés peuvent mesurer 6, 7, 14, 21 ou 42 mm de côté.
\item[E.] Il faut planter des piquets tous les 6 mètres.
\item[] Il faudra donc 94 piquets.
\item[F.] Ils se retrouveront dans les mêmes positions après 6\,140\,400 h (255\,850 jours $\approx$ 700,5 années).
\item[] 46\,200 révolutions pour Io, 72\,240 pour Europe, 35\,700 pour Ganymède et 15\,351 pour Callisto.
\item[G.] Le vigneron possède 900 bouteilles.
\item[I.] Henri possède CHF 432.--
\item[K.] Le prix est de CHF 6, 9, 12, 18 ou 36.--.
\item[] Avec un maximum de 120 personnes, le prix est de CHF 9.--.
\item[L.] 176 élèves participent au camp d'été.
\item[M.] Elle a perdu 5 membres.
\item[] La cotisation revient à CHF 15.--
\item[N.] Chaque enfant a reçu 49 billes chacun et il en reste 6 pour la maîtresse.
\end{itemize}

\vspace{20pt}
% =========================================================
\etape{4 \quad Puissances}
% =========================================================

\header{NO62 --- Puissances à calculer}
\vspace{-15pt}\begin{multicols}{4}
\begin{enumerate}[leftmargin=*,align=parleft, label=\alph*)]
\item 343
\item 144
\item 10\,000\,000
\item 1 
\item 1\,326
\item 128
\item 125
\item 1
\end{enumerate}
\end{multicols}

\newpage
\header{NO63 --- Puissances à ordonner}
\vspace{-15pt}\begin{enumerate}[leftmargin=*,align=parleft, label=\alph*)]
\item $2^7 > 2^6 > 2^5 > 2^4 > 2^3 > 2^2 > 2^1$
\item $7^5 > 6^5 > 5^5 > 4^5 > 3^5 > 1^5$
\item $10^3 > 3^3 > 3^2 > 2^3 > 1^{12}$
\item $10^5 > 3^4 > 2^6 > 5^2 > 4^2 > 2^3 > 4^1$
\end{enumerate}

\vspace{-5pt}\header{NO64 --- Au départ}
\vspace{-15pt}\begin{multicols}{4}
\begin{enumerate}[leftmargin=*,align=parleft, label=\alph*)]
\item 4
\item 8
\item 4
\item 10 
\item 100 et 50
\item 40 et 400
\item 100 et 10  
\item 3 et 9 
\item 400 et 100
\item 30 et 900
\end{enumerate}
\end{multicols}

\vspace{20pt}
% =========================================================
\etape{5 \quad Priorité des opérations}
% =========================================================

\header{NO78 --- Prioritaire}
\begin{multicols}{7}
\vspace{-15pt}\begin{enumerate}[leftmargin=*,align=parleft, label=\alph*)]
\item 14
\item 14
\item 16
\item 375
\item 4
\item 33
\item 49
\item 25
\item 27
\item 2
\item 1
\item 7
\item 12
\item 14
\item 7
\item 4
\item 16
\item 0
\item 51
\item 18
\item 4
\end{enumerate}
\end{multicols}

\vspace{-5pt}\header{NO82 --- Calcul réfléchi}
\vspace{-15pt}\begin{multicols}{3}
\textbf{a)}\\
2520\\
1,2\\
4,2\\
20\\
6,25\\
13,8\\
710\\
8\\
0,9\\
0,001

\textbf{b)}\\
 500\\
 2,5\\
 20,1\\
 70\\
 14\\
 24\\
 2,8\\
 0,01\\
 12,2\\
0,09

\textbf{c)}\\
19\\
260\\
210\\
11\\
12,6\\
7,2\\
0,004\\
0,32\\
36\,000\\
24\\
\end{multicols}

\vspace{-5pt}\header{NO100 --- De gauche à droite\,?}
\vspace{-15pt}\begin{multicols}{5}
\begin{enumerate}[leftmargin=*,align=parleft, label=\alph*)]
\item 24
\item 0,6
\item 31,5 
\item 22,5 
\item 7  
\item 34,5  
\item 14,9
\item 54,9 
\item 47,5 
\item 5,5
\end{enumerate}
\end{multicols}

\newpage
\header{NO89 --- Dans la vie courante}
\vspace{-10pt}\begin{enumerate}[leftmargin=*,align=parleft, label=\alph*)]
\item Je vais payer Fr.\,20.25.
\item Il lui faudra 4,4\,h (soit 4\,h\,24) pour recouvrir le sol de la salle de bain.
\item Il peut transporter 522 passagers.
\item Je dépense Fr.\,5.75.
\item Il paye Fr.\,600.-- au mètre, et seulement Fr.\,525. s'il peut négocier pour 3,5\,m
\item On doit lui rendre Fr.\,43.40.
\item Le père mesure 1,69\,m.
\end{enumerate}

\header{NO98 --- De tête\,!}
\vspace{-15pt}\begin{multicols}{5}
\begin{enumerate}[leftmargin=*,align=parleft, label=\alph*)]
\item 20,11
\item 42
\item 27,9
\item 72
\item 500
\item 24,7
\item 34
\item 1\,200
\item 25
\item 16
\item 20
\item 14
\item 66,5
\item 9,3
\item 351
\item 0,007
\item 150
\item 75
\item 111
\item 31
\end{enumerate}
\end{multicols}

\header{NO99 --- Oralement\,!}
\vspace{-15pt}\begin{multicols}{5}
\begin{enumerate}[leftmargin=*,align=parleft, label=\alph*)]
\item 1\,408
\item 0,09
\item 0,17
\item 0,03 
\item 16 
\item 750
\item 24
\item 360
\item 1,28
\item 5
\end{enumerate}
\end{multicols}

\header{NO75 --- Nombres croisés}
\vspace{-10pt}\begin{minipage}[t]{0.45\textwidth}
a)\vspace{2pt}

\begin{tabular}{c|ccc}
 & D & E & F \\
\hline
A & 3 & 2 & 5 \\
B & 2 & 4 &   \\
C & 2 & 0 & 7 \\
\end{tabular}
\end{minipage}%
\begin{minipage}[t]{0.45\textwidth}
b)\vspace{2pt}

\begin{tabular}{c|cccc}
 & E & F & G & H \\
\hline
A &   & 1 & 2 & 5 \\
B &   & 7 &   &   \\
C & 6 & 2 & 4 & 3 \\
D & 4 &   & 6 & 6 \\
\end{tabular}
\end{minipage}

\header{NO76 --- D'autres nombres croisés}
\vspace{-10pt}\begin{minipage}[t]{0.25\textwidth}
a)\vspace{2pt}

\begin{tabular}{c|cccc}
 & E & F & G & H \\
\hline
A & 6 & 5 & 6 & 1 \\
B & 7 & 8 &   & 3 \\
C &   & 6 & 2 & 5 \\
D & 9 & 4 & 4 & 7 \\
\end{tabular}
\end{minipage}%
\begin{minipage}[t]{0.25\textwidth}
b)\vspace{2pt}

\begin{tabular}{c|cccc}
 & E & F & G & H \\
\hline
A & 6 & 5 & 7 &   \\
B & 3 & 4 & 5 & 6 \\
C & 2 & 0 & 4 & 8 \\
D & 1 & 3 & 3 & 1 \\
\end{tabular}
\hfill
\end{minipage}%
\begin{minipage}[t]{0.25\textwidth}
c)\vspace{2pt}

\begin{tabular}{c|cccc}
 & E & F & G & H \\
\hline
A & 3 & 5 & 7 & 9 \\
B & 4 & 0 & 9 & 6 \\
C & 4 & 1 &   & 3 \\
D & 3 & 9 & 6 & 0 \\
\end{tabular}
\end{minipage}
\begin{minipage}[t]{0.25\textwidth}
d)\vspace{2pt}

\begin{tabular}{c|ccccc}
 & F & G & H & I & J \\
\hline
A & 7 & 2 &   & 6 & 4 \\
B & 3 & 4 & 5 & 6 & 7 \\
C &   & 6 & 2 & 5 &   \\
D & 8 & 4 & 0 & 6 & 2 \\
E & 9 &   & 9 & 6 & 1 \\
\end{tabular}
\end{minipage}

\newpage
% =========================================================
\etape{6 \quad Exercices \og\,Défi \fg}
% =========================================================

\header{NO43 --- Énergie solaire}
\vspace{-10pt}Il était alors 6\,h\,30.

\vspace{20pt}\header{NO73 --- À travers le monde}
\vspace{-5pt}\begin{tabularx}{\textwidth}{lCCCC}
Numération décimale & 8 & 73 & 336 & 1540 \\
Numération égyptienne   &  &  &  &  \\
Numération romaine   & VIII & LXXIII & CCCXXXVI & MDXL \\
Numération grecque   &  &  &  &  \\
Numération binaire   & 1000 & 1001001 & 101010000 & 11000000100 \\
\end{tabularx}

\vspace{20pt}\header{NO60 --- Cryptarithmes}
\vspace{-15pt}\begin{enumerate}[leftmargin=*,align=parleft, label=\alph*)]
\item Il y a quatre solutions:
\item[] 1735, 1735, 1735, 5205
\item[] 1745, 1745, 1745, 5235
\item[] 1765, 1765, 1765, 5295
\item[] 1795, 1795, 1795, 5385
\item Il y a deux solutions: $656+656=1312$ et $858+858=1716$
\item $E=0$.
\item[] Si $B=4$ (et $R=2$): il y a deux solutions ($M$ et $P$ valent 1 et 3).
\item[] Si $B=6$ (et $R=3$): il y a quatre solutions ($M$ et $P$ valent 1 et 5 ou 2 et 4).
\item[] Si $B=8$ (et $R=3$): il y a six solutions ($M$ et $P$ valent 1 et 7 ou 2 et 6 ou 3 et 5)
\item Il y a une seule solution: $43127+95127=138254$
\item C'est impossible.
\item Il y a deux solutions: $31840+578240=71840+538240=610080$
\end{enumerate}

\vspace{20pt}\header{NO15 --- Des félicitations à la clé}
\vspace{-10pt}Les cinq lettres sont V, A, R, B, O et peuvent former le mot BRAVO.

\vspace{20pt}\header{NO51 --- Toujours premier\,?}
\vspace{-10pt}Non, puisqu'elle n'est en tout cas pas valable pour 11 ni pour aucun multiple de 11.

\vspace{20pt}\header{Problème 2 --- Triplets Pythagoriciens}
\vspace{-10pt}\begin{minipage}[t]{0.48\textwidth}
A.\vspace{2pt}

\begin{tabular}{ccc}
X & Y & Z \\
\hline
3 & 4 & 5 \\
5 & 12 & 13 \\
7 & 24 & 25 \\
9 & 40 & 41 \\
11 & 60 & 61 \\
13 & 84 & 85 \\
15 & 112 & 113 \\
17 & 144 & 145 \\
\end{tabular}
\end{minipage}%
\begin{minipage}[t]{0.48\textwidth}
B.\vspace{2pt}

\begin{tabular}{ccc}
X & Y & Z \\
\hline
238 & 240 & 338 \\
\end{tabular}
\end{minipage}

